\documentclass[11pt]{amsart}
\usepackage{geometry}                % See geometry.pdf to learn the layout options. There are lots.
\geometry{letterpaper}                   % ... or a4paper or a5paper or ... 
%\geometry{landscape}                % Activate for for rotated page geometry
%\usepackage[parfill]{parskip}    % Activate to begin paragraphs with an empty line rather than an indent
\usepackage{graphicx}
\usepackage{amssymb}
\usepackage{epstopdf}
\DeclareGraphicsRule{.tif}{png}{.png}{`convert #1 `dirname #1`/`basename #1 .tif`.png}

\usepackage{enumerate}
\usepackage{nicefrac}

\usepackage{mathrsfs}

\newcommand{\R}{\mathbb{R}}
\newcommand{\D}{\mathbb{D}}
\newcommand{\tstD}{\mathscr{D}}
\newcommand{\semP}{\mathscr{P}}
\newcommand{\eps}{\varepsilon}
\newcommand{\supp}{\mathrm{supp}}
\newcommand{\exd}{\mathrm{d}}
\newcommand{\loc}{\mathrm{loc}}

\title{Math 580 Take home midterm exam 1}
\author{Due Wednesday October 23}
\date{Fall 2013}                                           % Activate to display a given date or no date

\begin{document}
\maketitle

\begin{enumerate}[1.]
\item
Let $\Omega$ be a domain, and let $\Sigma=\partial\Omega\cap B$ be a smooth and nonempty portion of the boundary, where $B$ is an open ball.
Let $u\in C^2(\Omega)\cap C^1(\Omega\cup\Sigma)$ satisfy $\Delta u=0$ in $\Omega$ and $u=\partial_\nu u = 0$ on $\Sigma$.
Show that $u$ is identically zero in $\Omega$.
\item
Consider the problem of minimizing the energy
$$
Q(u) = \int_I \left( 1+|u'(x)|^2 \right)^{\frac14} \exd x,
$$
for all $u\in C^1(I)\cap C(\overline I)$ satisfying $u(0)=0$ and $u(1)=1$, where $I=(0,1)$.
Show that the infimum of $Q$ over the admissible functions is $1$, but this value is not assumed by any admissible function.
\item
Consider the Dirichlet problem on the unit disk $\D\subset\R^2$ with the homogeneous Dirichlet boundary condition.
The solution is $u\equiv0$, which also minimizes the Dirichlet energy
$$
E(u) = \int_{\D}|\nabla u|^2.
$$
Construct a sequence $\{u_k\}\subset C(\overline\D)$ of functions satisfying all of the following conditions.
\begin{itemize}
\item
$u_k$ is piecewise smooth and $u_k|_{\partial\D}=0$ for all $k$,
\item
$E(u_k)\to0$ as $k\to\infty$, and
\item
$u_k$ as $k\to\infty$ diverges in a set that is dense in $\D$.
\end{itemize}
Then show that $u_k\to0$ in $H^1(\D)$.
\item
Exhibit a sequence $\{v_k\}\subset \tilde C^1(\D)$ that is Cauchy with respect to the $H^1$-norm,
whose limit is not essentially bounded.
\item
In this exercise we will study Sobolev spaces on the interval $I=(0,1)$.
Let $1\leq p < \infty$, and define the norm
$$
\|u\|_{1,p} = \left( \|u\|_{L^p}^p + \|u'\|_{L^p}^p \right)^{1/p},
$$
for $u\in C^1(\overline I)$.
Let $H^{1,p}(I)$ be the completion of $C^1(\overline I)$ with respect to the norm $\|\cdot\|_{1,p}$.
\begin{enumerate}[a)]
\item
Show that there is a continuous injection of $H^{1,p}(I)$ into $L^p(I)$.
\item
Prove the {\em Sobolev inequality}
$$
\|u\|_{L^\infty} \leq 2^{1-1/p} \|u\|_{1,p},
\qquad u\in H^{1,p}(I).
$$
\item
Since $H^{1,p}(I)$ is a subspace of $L^p(I)$, an element of $H^{1,p}(I)$ is an equivalence class of functions that differ on sets of measure zero.
So one can change the values of a function on a set of measure zero, and it would still correspond to the same element in $H^{1,p}(I)$.
Make sense of, and prove the statement that the elements of $H^{1,p}(I)$ are continuous functions.
\item
Prove the {\em Friedrichs inequality}
$$
\|u\|_{L^p}^p \leq 2^{p-1} \|u'\|_{L^p}^p + 2^{p-1} |u(\xi)|^p,
\qquad u\in H^{1,p}(I),
\quad \xi\in[0,1].
$$
In particular, make sense of the derivative $u'$ appearing in the right hand side.
\item
Let $H^{1,p}_0(I)$ be the closure of $C^1_c(I)$ in $H^{1,p}(I)$.
Show that 
$$
H^{1,p}_0(I) = \{u\in H^{1,p}(I) : u(0)=u(1)=0\}.
$$
\item
Prove the {\em Poincar\'e inequality}
$$
\|u\|_{L^p}^p \leq 2^{p-1} \|u'\|_{L^p}^p + 2^{p-1} \big| \int_I u \big|^p,
\qquad u\in H^{1,p}(I).
$$
\end{enumerate}
\item
\begin{enumerate}[a)]
\item
Let $\Omega\subset\R^n$ be a bounded domain and let $g\in H^1(\Omega)$.
Give a detailed proof of the fact that there exists a unique $u\in H^1(\Omega)$ satisfying $u-g\in H^1_0(\Omega)$ and
$$
\int_\Omega \nabla u \cdot \nabla v = 0
\qquad\textrm{for all}\quad 
v\in\tstD(\Omega).
$$
\item
In the context of a),
let us denote the map that sends $g$ to $u$ by 
$$
S:g\mapsto u:H^1(\Omega)\to H^1(\Omega).
$$
Show that $S$ is linear and is Lipschitz continuous in the sense that there is a constant $c$ such that
$$
\|S(g_1)-S(g_2)\|_{H^1} \leq c \|g_1-g_2\|_{H^1},
\qquad
g_1,g_2\in H^1(\Omega).
$$
\item
Exhibit an example of an unbounded domain $\Omega$ and a sequence $\{u_k\}\subset H^1_0(\Omega)$ 
such that $\{u_k\}$ minimizes the Dirichlet energy on $\Omega$ and that it is not Cauchy with respect to the $H^1$-norm.
\end{enumerate}
\end{enumerate}

\end{document}  












